\section{Conclusion} \label{sec:conclusion}
This paper proposes an improved GoP for dysarthria speech intelligibility assessment by using UQ methods.
Expected to alleviate the problem of modern NN's overconfidence, especially for disordered speech, tested UQ methods include (1) normalization of phoneme prediction and (2) modification of the scoring function. 
The experiments were carried out on dysarthric speech datasets in English, Korean, and Tamil.
According to the experimental results, the prior normalized MaxLogit GoP shows the best performance, outperforming both the traditional GoPs and other proposed GoP variants.
Furthermore, to verify the usefulness of our proposed method, an analysis of which phoneme pronunciation scores highly correlate to speech intelligibility is conducted.


\section{Acknowledgements}
This work was supported by Institute of Information \& communications Technology Planning \& Evaluation (IITP) grant funded by the Korea government (MSIT) (No.2022-0-00223, Development of digital therapeutics to improve communication ability of autism spectrum disorder patients).

% A limitation of this study is the requirement for time alignment, which will be addressed in future research. 
% Additionally, other speech dimensions such as prosody must also be taken into consideration for effective speech intelligibility assessment.
% This paper contributes in that it introduces an improved GoP, considering the 


% This paper contributes in that it showcases the effectiveness of GoP, which has previously been underutilized in the study of disordered speech.
% Its ability to provide interpretable results at the phoneme level without the parallel corpus could prove valuable in the clinical field.
% % Secondly, the effectiveness of UQ methods on overconfident predictions from neural networks has been investigated.
% % Experiments verified that UQ methods can address the problem, especially in the disordered speech domain.
% % This research contributes to automatic speech intelligibility assessment by demonstrating the potential use of bot
% % The findings suggest that significant acoustic differences must be taken into account when using probabilities from neural networks, especially in applications to disordered speech.
% A limitation of this study is the requirement for time alignment, which will be addressed in future research. 
% Additionally, other speech dimensions such as prosody must also be taken into consideration for effective speech intelligibility assessment.




% The analysis also reveals that while different phonemes have varying relationships with speech intelligibility across different languages, fricatives tend to have a higher correlation with speech intelligibility.
% The experiments conducted on English, Korean, and Tamil dysarthric speech datasets, results indicate that the prior normalized MaxLogit GoP variant is the most effective. 
% This paper contributes to the field of automatic speech intelligibility assessment, by showcasing the application of improved GoP for disordered speech, especially dysarhtric speech.
% Experimental results imply that significant acoustic differences between dysarthric and healthy must be considered, especially when probabilities from the neural networks are employed.

% This paper contributes in showcasing the effectiveness of UQ methods to alleviate the significant differences between dysarthric and healthy speech.

% The study takes into account the significant differences between dysarthric speech and healthy speech.

% This paper contributes in showcasing the effectiveness of GoP, which has previously been underutilized in the study of disordered speech.
% Its ability to provide interpretable results at the phoneme level without the parallel corpus could prove valuable in the clinical field.
% Moreover, the effectiveness of UQ methods on uncalibrated results from the neural networks has been investigated.
% Experiments verified that UQ methods can address the over-confident and uncalibrated outputs produced by neural networks in the speech domain.

%Limitation

% This paper makes two contributions to the field of automatic speech intelligibility assessment. 
% Firstly, it showcases the effectiveness of GoP, which has previously been underutilized in the study of disordered speech.
% Its ability to provide interpretable results at the phoneme level without the parallel corpus could prove valuable in the clinical field.
% Secondly, the effectiveness of UQ methods on uncalibrated results from the neural networks has been investigated.
% Experiments verified that UQ methods can address the over-confident and uncalibrated outputs produced by neural networks in the speech domain.
% Limitation


% This paper investigates the use of UQ methods to enhance the accuracy of GoP variants for assessing dysarthric speech intelligibility. 
% The study takes into account the significant differences between dysarthric speech and healthy speech.
% Applied UQ methods include (1) scoring function modification and (2) phoneme prediction modification.
% Tested on TORGO English, QoLT Korean, and SSNCE Tamil dysarthric speech datasets, the best GoP variant is found to be prior normalized MaxLogit GoP.
% Further analysis of each phoneme implies while phonemes have different relations to speech intellgiblity by languages, fricative \textipa{/s/} shared the pattern to have high correlations to speech intelligibility.

% While the overall results of the experiments using UQ methods outperformed the baseline experiments, Normalized MelLogit-Gop showed the greatest correlation with the intelligibility scores.
% Experimental results suggest that traditional GoPs were not properly calibrated, and UQ methods can effectively address this issue.
% Furthermore, we conduct phoneme-level analysis to find frequently mispronounced phonemes by each language.
% This paper makes two contributions to the field of automatic speech intelligibility assessment. 
% Firstly, it showcases the effectiveness of GoP, which has previously been underutilized in the study of disordered speech.
% Its ability to provide interpretable results at the phoneme level without the parallel corpus could prove valuable in the clinical field.
% Secondly, the effectiveness of UQ methods on uncalibrated results from the neural networks has been investigated.
% Experiments verified that UQ methods can address the over-confident and uncalibrated outputs produced by neural networks in the speech domain.
%Limitation

% an improved version UQ method on GoP, and demonstrates its effectiveness 

% A limitation of our study lies in the discrepancy between the evaluation and labeling methods.
% While the labels were assigned based on a speaker level, we assessed the intelligibility scores on an utterance level, assuming that all utterances from the speaker had the same level of intelligibility.
% To alleviate such issues, we plan to aggregate the intelligibility scores of individual utterances into speaker-level scores.
% Another limitation is that we did not provide information on how each phoneme was mispronounced. 
% Hence, we plan to develop a model that can provide this level of detail in our future work.


%% Consider Percpetual evaluation + other aspects other than pronunciation.

% , which will help us better understand the specific areas where non-native speakers struggle with pronunciation.
% no need of large trained models





% First, the analysis is limited to sentences. 
% Future work includes 



% In addition, cross-lingual representations from the wav2vec2 XLS-R model are extracted, to overcome the limitations of most GoPs where language-specific acoustic models with large amounts of training data is required.

% The GoP with UQ method is tested on dysarthric datasets in English, Korean, and Tamil, and results in higher Kendall's rank coefficients with speech intelligibility classes than previous experiments using GMM-GoP, NN-GoP, ASR free-GoP, and SSL-GoP.




